\section{Open Questions}

The following five questions were not answered by this replication and
represent concrete, tractable follow-ups.

\subsection*{Q1. Extension to time-dependent Hamiltonians}
\textbf{Question:} Do the empirical PF1/PF2/PF4 error-scaling exponents
survive extension from time-independent to time-dependent Hamiltonians
$H(t)$ (e.g.\ driven Heisenberg, Floquet models), and do the
Suzuki--Magnus product formulas recover the same $p+1$ order?

\textbf{Basis:} The replication and the paper both restrict to
time-independent $H$. Real quantum-simulation targets (driven
condensed-matter, molecular dynamics, quantum control) are often
time-dependent, where the Magnus expansion introduces additional
commutator terms and the empirical-vs-bound gap may narrow or invert.

\textbf{Next steps:} Extend the $n=6$ Heisenberg benchmark by turning on
a time-dependent transverse field $h_x(t)=A\cos(\omega t)$. Implement
standard Magnus-based PF2/PF4 variants (e.g.\ commutator-free 4th-order
Magnus). Recompute the log--log slope of
$\|U_{\text{exact}}(t)-U_{\text{MagnusPF}}(t)\|_2$ vs $r$ for several
$(A,\omega)$. Compare fitted slopes against theoretical $p+1$ for
commutator-free schemes.

\subsection*{Q2. Hybrid LCU / qubitization + product-formula scheme}
\textbf{Question:} How does a hybrid LCU / qubitization + product-formula
scheme (Trotterize the low-norm block, qubitize the high-norm block)
compare in gate count against pure PF4 or pure QSP on the paper's
specific Heisenberg + random-field Hamiltonian at accessible sizes
$n=8$--$14$?

\textbf{Basis:} The paper compares PF, LCU, and QSP as monolithic
algorithms. Hybrid PF+LCU decompositions (Low--Wiebe, Childs--Wiebe
splitting) have not been quantitatively benchmarked against the paper's
own Hamiltonian class at empirically-scaled resource cost.

\textbf{Next steps:} Split $H$ into $H_{\text{bulk}}$ (XX+YY+ZZ bonds)
and $H_{\text{field}}$ (single-site random $Z$). Implement (i) pure PF4
(baseline), (ii) LCU on $H_{\text{field}}$ + PF4 on $H_{\text{bulk}}$
interleaved, (iii) pure qubitization on the full block-encoding. Report
empirical error vs total controlled-gate count on $n\in\{8,10,12,14\}$
at target error $10^{-4}$.

\subsection*{Q3. Chemistry / lattice benchmarks}
\textbf{Question:} For a concrete chemistry benchmark
(H\textsubscript{2}, LiH, H\textsubscript{2}O in a minimal basis) or a
lattice benchmark (2D Fermi--Hubbard $4\times 4$ at half filling), does
the PF4 empirical-vs-bound gap remain the ${\sim}4\times$ factor
observed for random-field Heisenberg, or does it grow / shrink materially
with Hamiltonian structure?

\textbf{Basis:} The paper focuses on spin systems; its empirical-vs-bound
gap is measured on Heisenberg. Chemistry Hamiltonians have very different
term-structure, so the effective empirical constant may differ by orders
of magnitude, changing the resource-estimate takeaway.

\textbf{Next steps:} Use OpenFermion + PySCF (free) to generate
H\textsubscript{2}/LiH/H\textsubscript{2}O JW-transformed Hamiltonians in
STO-3G. Split into a fermionic-swap PF4 partitioning. Measure
spectral-norm empirical error vs $r$ and compare against the analytic
commutator bound. Repeat for $4\times 4$ Fermi--Hubbard on a
laptop-tractable subsector. Report $\text{ratio}(r)$ trajectory.

\subsection*{Q4. Robustness under realistic gate noise}
\textbf{Question:} How robust is the empirical PF4 error advantage under
realistic gate-level noise (depolarizing plus coherent overrotation) at
fault-tolerant-relevant physical-error rates ($10^{-3}$ to $10^{-4}$)?
Does the $r^{-4}$ scaling survive, or collapse to a noise-floor plateau
at moderate $r$?

\textbf{Basis:} The paper's error model is exact algorithmic Trotter
error; it does not address noise--Trotter interference. In a real
fault-tolerant setting, logical-gate errors will eventually dominate
algorithmic error, potentially eliminating the PF4 advantage at large $r$.

\textbf{Next steps:} Wrap each 2-qubit gate in the PF4 circuit with a
Kraus-map depolarizing channel at $p\in\{10^{-3},10^{-4},10^{-5}\}$.
Compute average-gate-infidelity of noisy $U_{\text{PF4}}$ vs
$U_{\text{exact}}$ via density-matrix simulation on $n=4$--$6$. Plot
infidelity vs $r$ for each $p$ and identify the crossover $r^\star$
where noise dominates. Compare crossover locations for PF2 vs PF4.

\subsection*{Q5. Sample-cost Pareto frontier}
\textbf{Question:} What is the true sample-cost / resource Pareto
frontier between PF1, PF2, PF4 (and a bespoke nearly-optimal Trotter
product formula) on the paper's Hamiltonian at fixed error target
$\varepsilon\in\{10^{-2},10^{-4},10^{-6}\}$, once realistic Clifford+T
compilation is folded in?

\textbf{Basis:} The scaling-exponent comparison shows PF4 always wins
asymptotically, but at practical $\varepsilon$ the constant-factor
per-step gate count of PF4 (many nested $S_2$ layers) can make PF2
cheaper. A full Pareto sweep over $\varepsilon$ is not provided by the
paper.

\textbf{Next steps:} Compile each PF1/PF2/PF4 step to Clifford+T using a
free toolchain (e.g.\ \texttt{staq} or \texttt{PyZX}). Count total
T-gates for the $r$ required to reach $\varepsilon\in\{10^{-2},10^{-4},10^{-6}\}$
on $n=6$ Heisenberg. Plot total T-count vs $\varepsilon$ for each order.
Overlay the paper's asymptotic scaling to identify the $\varepsilon$
range in which PF4 actually dominates PF2 in constant-factor terms.
